\section{Discusión}

\subsection{6.1 Ventajas e Implicaciones}

Resulta revelador observar cómo un mejor equilibrio entre planificación y control puede transformar radicalmente la viabilidad de infraestructuras orbitales de gran escala. Los resultados obtenidos muestran reducciones consistentes en tiempo de ensamblaje y consumo energético, junto con una disminución notable en el índice de riesgo compuesto. Estas mejoras no son meramente incrementales. 

Para misiones de exploración profunda, donde la ventana de oportunidad es estrecha y la logística extremadamente costosa, disponer de sistemas capaces de ensamblar telescopios o habitats modulares de forma autónoma cambia el paradigma. La coordinación multi-brazo y el uso intensivo del gemelo digital permiten responder a imprevistos sin intervención terrestre, un aspecto clave cuando las comunicaciones sufren demoras de minutos u horas. 

\begin{table}[h]
\centering
\caption{Brechas identificadas y grado de mitigación logrado por el framework propuesto}
\label{tab:brechas_identificadas}
\begin{tabular}{lcc}
\hline
\textbf{Brecha} & \textbf{Estado del Arte} & \textbf{Framework Propuesto} \\
\hline
Eficiencia temporal & Media & Alta \\
Robustez ante perturbaciones & Baja-Media & Alta \\
Autonomía en exploración profunda & Limitada & Mejorada \\
Consumo energético & Elevado & Reducido 22\% \\
Escalabilidad multi-robot & Emergente & Demostrada en simulación \\
\hline
\end{tabular}
\end{table}

Las implicaciones van más allá de la ingeniería. Facilitan arquitecturas espaciales reparables y evolutivas, reduciendo la dependencia de lanzamientos únicos y costosos.

\subsection{6.2 Limitaciones y Desafíos Abiertos}

Uno de los desafíos persistentes que surge al analizar sistemas de este tipo radica en la transición desde simulación a hardware real. Aunque los resultados en entorno controlado son sólidos, factores como degradación térmica de actuadores, variaciones en propiedades de materiales bajo radiación y fallos impredecibles de sensores siguen representando incertidumbres significativas.

El cómputo embarcado necesario para el gemelo digital y la replanificación en tiempo real impone restricciones de potencia y disipación térmica que no siempre serán fáciles de satisfacer en plataformas pequeñas. Además, la fidelidad del modelo dinámico tiende a degradarse cuando se incorporan estructuras flexibles de gran tamaño. No obstante, cabe matizar que estas limitaciones no invalidan el enfoque, sino que señalan prioridades claras para iteraciones futuras.

\subsection{6.3 Comparación con Estado del Arte}

Lejos de ser un asunto resuelto, la literatura reciente sugiere que el campo avanza hacia mayor autonomía, aunque con ritmos desiguales. Revisiones exhaustivas de tecnologías de ensamblaje robótico in-orbit destacan precisamente la necesidad de marcos integrados como el aquí propuesto \cite{Wang2025_OASR}. 

De forma similar, análisis específicos sobre manipuladores autónomos para servicing y remoción de desechos enfatizan las dificultades de coordinación bajo incertidumbre \cite{Fallahiarezoodar2025_Manipulators}. Nuestra aproximación logra mejores compromisos entre velocidad y seguridad que muchos de los sistemas descritos en esas revisiones, especialmente en escenarios multi-componente.

La misión OSAM-1 de NASA, aunque finalmente cancelada, estableció un referente ambicioso en capacidades de ensamblaje y refuel \cite{NASA_OSAM1_2023}. El framework desarrollado aquí complementa aquellos esfuerzos al priorizar eficiencia y reducción cuantitativa de riesgos, elementos que resultan decisivos para operaciones sostenidas en el espacio profundo. 

En resumen, el trabajo aporta avances concretos en áreas donde el estado del arte todavía muestra debilidades estructurales. Queda camino por recorrer, particularmente en validación orbital real, pero los cimientos parecen sólidos.